\begin{textsample}{POS Dim 1 – ce – Score 67.00 – ce/ce\_em\_51.txt}  \label{ex:f1_pos_043}
8.4.7 A diversidade e a multiplicidade no ensino de Sociologia : narrativas, usos e formas de construção do conhecimento. 631 Pensar sobre o ensino da Sociologia no Ensino Médio \textbf{exige} principalmente \textbf{uma} reflexão constante sobre o sentido do aprendizado da Sociologia.
Diante disso, surgem algumas reflexões \textbf{que} cada professor da ciência deverá fazer e buscar responder no decorrer do seu processo de ensino aprendizagem : o \textbf{que} se quer \textbf{que} o aluno saiba sobre a Sociologia ou para \textbf{que} serve esta \textbf{disciplina} no ensino médio?
Como se organiza o ensino sociológico na escola para desenvolvimento dos indivíduos? 632 É no conhecimento das ciências humanas \textbf{que} torna -se possível encontrar os elementos epistemológicos \textbf{que} refletem sobre a realidade social.
Isto não quer \textbf{dizer} \textbf{que} essa condição não esteja presente nos demais conhecimentos, mas \textbf{que} essa é \textbf{uma} condição sine qua non desses saberes, em especial a Sociologia, \textbf{que} \textbf{historicamente} tem sua gênese relacionada ao surgimento da sociedade \textbf{moderna}. 633 Em sua gênese, a Sociologia recebe \textbf{forte} \textbf{influência} do pensamento positivista, se constituindo em conhecimento \textbf{que} se presta a explicar a nova ordem social capitalista.
Esta nova ciência, herdeira das \textbf{teorias} sociais do “ \textbf{século} das luzes ”, contribuiu para promover o entendimento e a aceitação de \textbf{uma} realidade social marcada pela concentração de riquezas por parte de \textbf{uma} minoria e geradora de pobreza e miséria social para a grande maioria. 634 A ideologia subjacente à nova ciência suscitou \textbf{que}, em determinado estágio da evolução das sociedades, os benefícios técnicos e materiais seriam acessíveis para todos os segmentos sociais, bem como indicaria o caminho a ser seguido por essa sociedade rumo à manutenção da ordem social para se alcançar o progresso da \textbf{humanidade}.
A reformulação \textbf{teórica} da Sociologia se deu a partir das reflexões de Karl Marx, \textbf{que} elaborou \textbf{uma} interpretação da realidade social constituída de \textbf{um} componente material, entendendo as relações sociais como se desenvolvendo a partir do modo como o ser humano busca satisfazer suas necessidades, situados no tempo, por isso, histórica. \textbf{Um} outro componente apontado por Marx foi o caráter \textbf{dialético} das relações, \textbf{uma} vez \textbf{que} elas trazem em si antagonismos e conflitos gerados pela luta entre as classes sociais.
Com isso, ele construiu \textbf{um} modelo \textbf{teórico} \textbf{baseado} na tese de \textbf{que} as relações na sociedade são materialmente determinadas ( determinismo econômico ).
A contribuição mais relevante de Marx ao conhecimento sociológico dá -se através do \textbf{método} denominado por ele de materialismo histórico \textbf{dialético}.
O \textbf{método} \textbf{dialético} e as correntes \textbf{que} o seguem servirão como \textbf{um} instrumento de crítica social, dando maior flexibilidade à Sociologia, fazendo com \textbf{que} esta deixe de ser \textbf{um} conhecimento em defesa da conservação da estrutura social e colocando -a em posição de suscitar mudanças e transformações na sociedade.
Contribuiu também para \textbf{uma} compreensão \textbf{dialética} da sociedade \textbf{contemporânea}, \textbf{revelando} sua \textbf{marca} capitalista e burguesa.
Viana ( 2006, p 15-18 ) define como sendo o longo processo de mudanças sociais \textbf{desencadeadas} no interior da sociedade feudal e \textbf{que} culminou com o estabelecimento de \textbf{uma} nova base de relações sociais definidas a partir do capital, elevando a burguesia à posição de classe. 635 O materialismo histórico acentuou o caráter inexorável da mudança da estrutura social e explicou como este fenômeno está \textbf{imbricado} na superação de \textbf{um} modo de produção \textbf{historicamente} determinado por outro.
A adoção do \textbf{método} \textbf{dialético} por parte dos herdeiros do marxismo tem repercutido em várias áreas do conhecimento, principalmente nas \textbf{Humanidades}.
As \textbf{teorias} críticas, no campo da educação, bem como as concepções pedagógicas progressistas advogam \textbf{que} só é possível a formação de \textbf{um} cidadão crítico e consciente sob a égide da concepção \textbf{dialética} da realidade. 636 Afirmar a Sociologia como conhecimento necessário à formação do aluno significa amplificar os referenciais epistemológicos capazes de promover mudanças na forma de pensar e agir da/o aluna/o, tornando -o agente questionador e transformador de sua realidade social.
Segundo ( Morais, 2004, p. 101 ) “ isso \textbf{revela} a competência permanente da Filosofia e da Sociologia, quer para formular as questões \textbf{que} interessam, quer para apresentar as respostas pertinentes.
Estão aí duas razões pelas quais não pode estar ausente do Ensino Médio ” ( p. 101 ) ”.
Neste sentido, o ensino de Sociologia permite ao \textbf{educando} a oportunidade de descobrir novos olhares sobre a realidade social com \textbf{que} este se defronta a partir dos saberes das Ciências Sociais quando esses conhecimentos possibilitam diálogos com as juventudes \textbf{que} se encontram na escola de ensino médio no Brasil. 637 A inclusão da Sociologia no currículo da Educação Básica no Brasil tem \textbf{vivido} alternâncias, \textbf{ora} pertencendo à base comum e sendo considerada \textbf{disciplina} obrigatória ; \textbf{ora} excluída do currículo obrigatório e se fazendo presente apenas como conteúdo transversal, utilizando parte dos 25 \% de carga horária destinada aos temas transversais.
Outras vezes, tem aparecido diluída entre os demais conhecimentos das Ciências Humanas.
As incertezas geradas por essa situação criam \textbf{um} clima de insegurança e dúvidas \textbf{que} atingem diretamente o trabalho com a \textbf{disciplina}. 638 \textbf{Um} retrocesso significativo ocorreu com o veto do ex-presidente Fernando Henrique Cardoso sobre a proposta do Deputado Federal Pe.
Roque \textbf{que} propunha incluir Sociologia como \textbf{disciplina} da base comum.
A supressão \textbf{que} o veto presidencial provocou contribuiu para \textbf{relegar} e distanciar, cada vez mais, esse conhecimento do processo de formação dos alunas/os, e, principalmente, da construção de \textbf{uma} educação \textbf{comprometida} com a democracia. ( Silva, 2004, p. 80 ) se posicionou a respeito dessa situação afirmando \textbf{que} “ esse ir e vir da Sociologia no Ensino Médio impediu \textbf{que} se desenvolvesse \textbf{uma} tradição de ensino desta ciência na escola.
Ficou prejudicada a pesquisa nesta área, o desenvolvimento de metodologias adequadas, de textos didáticos sérios, de recursos didáticos tais como audiovisuais, entre outros ” ( p. 80 ). 639 \textbf{Talvez} essa situação à qual a Sociologia tem sido \textbf{relegada} na Educação Básica, deva -se ao risco indicado por Bourdieu citado por ( Tomazi, 2004 ) “ a Sociologia não valeria nem \textbf{uma} hora de esforço se fosse \textbf{um} saber de especialista reservado aos especialistas ”.
Fortalecer a Sociologia no Ensino Médio \textbf{exige} inovações capazes de promover aprendizagens contextualizadas, \textbf{que} permitam aos alunas/os pensar sociologicamente.
A escola precisa tornar -se \textbf{um} espaço de síntese capaz de conduzir o aluno para \textbf{um} diálogo permanente sobre essa realidade contraditória, \textbf{que} \textbf{necessita} ser mediada pela compreensão e suscita propostas de mudanças sociais. 640 O ensino da Sociologia deve considerar aspectos relevantes à compreensão da \textbf{contemporaneidade}.
Nesta perspectiva, destacam -se o impacto e as modificações socioculturais produzidas pelas novas Tecnologias da Informação e Comunicação ( TIC ), como informadoras e formadoras de opiniões e como dimensão \textbf{central} das sociedades do \textbf{século} XXI.
É impossível ignorar o avanço e a \textbf{influência} das TIC no campo da educação, porquanto elas carregam em si \textbf{uma} atratividade \textbf{que} a escola, de \textbf{um} modo geral, não tem \textbf{conseguido} garantir. 641 No ensino da Sociologia, os recursos tecnológicos devem ser pensados e analisados em suas dimensões pedagógicas e ideológicas como instrumentos \textbf{que} podem elevar o nível cognitivo do aluno e como símbolo da sociedade pós-moderna, \textbf{que} busca nessas novas tecnologias, reorganizar o trabalho e o trabalhador.
A prática do professor deve propor \textbf{uma} ressignificação dessas tecnologias \textbf{que} devem estar a serviço de \textbf{uma} mudança na concepção pedagógica.
Deste modo, \textbf{talvez} se possa evitar \textbf{uma} prática de ensino tradicional e instrucionista supostamente renovadora, por conta da adoção destas novas tecnologias, mas conservadora, caso prepondere seu caráter tecnicista e instrumental. 642 Desenvolver \textbf{uma} prática de ensino \textbf{que} prepare o aluno a pensar com \textbf{categorias} sociológicas, \textbf{que} possibilite \textbf{uma} mediação da relação dele com a sociedade em \textbf{que} \textbf{vive} com o objetivo de desaliená -lo das ideologias vigentes, preparando o indivíduo para a cidadania, \textbf{revela} -se \textbf{uma} \textbf{tarefa} \textbf{bastante} árdua.
Esse processo de desalienação requer do aluno o exercício de \textbf{uma} crítica radical à sociedade e mudanças nas suas relações sociais, incluindo as mediadas pelas novas tecnologias. 643 Essa forma de ensinar Sociologia a distância daquele saber identificado com a ideia de \textbf{que} ele é especializado, \textbf{aproximando} -a da proposta de Bourdieu de torná -la acessível a todos.
Nas situações em \textbf{que} o ensino de Sociologia adotou a perspectiva de \textbf{um} saber eminentemente especializado, acabou provocando \textbf{um} distanciamento da prática epistemológica das alunas/os.
Portanto, afirmar a necessidade desse conhecimento à formação do aluno \textbf{implica} não só na sua inclusão na base comum, bem como no desenvolvimento de \textbf{uma} política educacional comprometida com \textbf{uma} educação capaz de proporcionar \textbf{uma} formação realmente crítica e reflexiva.
Reconhecer este fato \textbf{remete} à valorização do profissional de Sociologia, preparado intelectual e pedagogicamente para o ensino adequado dessa \textbf{disciplina}. 644 \textbf{Uma} pesquisa promovida pela Secretaria de Educação Básica do Estado do Ceará em 2003 \textbf{revelou} alguns dados importantes sobre o currículo de Sociologia do Ensino Médio.
Dentre eles, cabe destacar a evidente \textbf{fragmentação}, bem como a repetição dos conteúdos trabalhados de forma descontextualizada.
Além disso, muitos professores não têm formação específica e lecionam a \textbf{disciplina} apenas para completar a carga horária \textbf{que}, por lei, ele precisa dedicar à \textbf{docência}.
Tal constatação reforça ainda mais a necessidade de se definir \textbf{que} a Sociologia, bem como a Filosofia devem compor \textbf{disciplinas} da base comum.
As Diretrizes Curriculares para os cursos de graduação em Ciências Sociais estabelecem como competências e habilidades do profissional o “ domínio dos conteúdos básicos \textbf{que} são objetos de ensino e aprendizagem nos Ensinos Fundamental e Médio ” ( MEC, 2001 ). 645 O caminho a ser trilhado pelo professor para a reconstrução e transformação de sua práxis no ensino da Sociologia, passa por \textbf{um} processo reflexivo, podendo vir a constituir -se num instrumento efetivo de mudança, através do questionamento da realidade. \textbf{Um} aspecto delicado, mas \textbf{que} não pode ser desconsiderado, se refere à postura ideológica do professor \textbf{que} pode ser assumida ou não.
Neste sentido, faz -se necessário \textbf{que}, para além das suas preferências, ele não sacrifique conteúdos e muito menos assuma \textbf{uma} postura doutrinária, comprometendo o aprendizado do aluno. 646 A escola e a educação se tornam instrumentos \textbf{que} ajudam a legitimar qualquer ordem de natureza ideológica, pelo papel social \textbf{que} \textbf{desempenham} junto à sociedade - de instrumentalização cultural a partir do conhecimento das ciências.
As formas e os modos como essa instrumentalização acontece constitui a problemática \textbf{central} na discussão sobre o currículo, pois ao invés de servir de referência para \textbf{que} o \textbf{sujeito} possa ir além do olhar \textbf{que} já possui sobre o mundo do qual participa, ele pode indicar \textbf{uma} possibilidade de aprisionamento onde o conformismo social torna -se num \textbf{imperativo} cultural. 647 Florestan Fernandes ( 1963 ) \textbf{um} dos precursores da Sociologia brasileira e defensor da inclusão dessa ciência como \textbf{disciplina} obrigatória no ensino secundário no Brasil, em análise sobre o papel da Sociologia no mundo \textbf{moderno} sinaliza \textbf{que} seus intelectuais precisam conhecer de modo profundo as verdadeiras intenções dos grupos \textbf{que} detém o poder econômico e político, pois os saberes sociológicos podem servir de instrumentos de repressão e controle sociais quando utilizados pela classe dominante. 648 Nenhuma proposta curricular está isenta de intenções.
Nesse aspecto, o olhar do professor é determinante na consecução de seus propósitos.
A autonomia pode servir para práticas sociais distintas, \textbf{ora} criando possibilidades de leitura e intervenção críticas pelos alunas/os sobre o mundo e a realidade vivenciada por estes, \textbf{ora} colaborando para ações \textbf{que} visam tão somente o conformismo e a conveniência de quem produziu.
Produzir \textbf{uma} proposta de currículo pressupõe \textbf{uma} atitude política para a ação pedagógica no ensino. 649 A ensino da Sociologia no Ensino Médio tem representado \textbf{um} \textbf{expressivo} ganho na formação das novas gerações \textbf{que} frequentam a escola pública.
É importante ressaltar a participação \textbf{que} os professores, em especial os de Sociologia, têm na elaboração de \textbf{uma} proposta curricular \textbf{que} proponha \textbf{uma} escola \textbf{comprometida} na formação e compreensão sociológica da realidade social, o \textbf{que} \textbf{exige} a efetiva presença dessa \textbf{disciplina} no currículo. 650 A Sociologia na escola pública é necessária, importante e deve ser reconhecida pelos \textbf{sujeitos} envolvidos na construção do projeto pedagógico da escola, incluindo os gestores, o corpo docente, \textbf{discente}, funcionários e a comunidade.
No caso do Ceará, a inclusão da Sociologia como \textbf{disciplina} no Ensino Médio, já vem de alguns anos, e consta nas Diretrizes para Educação Básica desde 2004.
Efetivamente se estabeleceu como \textbf{disciplina} obrigatória no currículo, na era do governo Lula, perdurando até 2017, quando \textbf{perdeu} a obrigatoriedade com a reforma do ensino médio.
Todavia, no Estado do Ceará, a proposta do ensino da Sociologia, assim como da Filosofia, têm se mantido no currículo. 651 Grande parte do trabalho docente e processo pedagógico, também com o ensino da Sociologia, desenvolve -se em sala de aula, portanto, o professor deve utilizar sua sensibilidade e criatividade para tornar este espaço mais vivo e dinâmico.
Para tanto, deve se munir de variadas estratégias, rompendo o predomínio da aula expositiva, \textbf{que} embora tenha seu valor, não deve ser a única estratégia pedagógica.
Dentre variadas estratégias, pode -se promover a discussão de \textbf{um} problema social previamente estudado ou organizar a classe em equipes para, a partir de subtemas diferenciados, aprofundar \textbf{um} determinado tema.
É possível também estimular a produção de pequenos textos, individualmente ou em grupos, objetivando desenvolver a capacidade de elaboração de análises \textbf{que} exercitem o \textbf{raciocínio} lógico, reflexivo e crítico. 653 Sugere -se a criação de mecanismos \textbf{que} exercitem a imaginação sociológica, levando o aluno a distinguir problemas sociais de problemas sociológicos, fazendo -o também perceber a diferença entre senso comum e conhecimento científico.
Deve -se levar o aluno, efetivamente, a elaborar pesquisas científicas de forma individual ou em grupo, contribuindo para \textbf{que} ele compreenda \textbf{que} pesquisar pressupõe problematizar, compreender a realidade, buscando se situar no mundo \textbf{contemporâneo}. 654 É consenso entre os educadores \textbf{que} o processo pedagógico pode se tornar muito profícuo, à medida \textbf{que} extrapola a sala de aula.
Portanto, se a escola dispõe de sala de vídeo, biblioteca, laboratório de informática, Centro de Multimeios, ou outros espaços pedagógicos diferenciados, o professor pode utilizar, da forma mais criteriosa possível, estes espaços, possibilitando aos seus alunas/os \textbf{estímulos} cognitivos diferenciados, além de trabalhos de campo, visitas culturais e outros \textbf{que} saem dos muros e paredes da escola, onde também se faz a construção do conhecimento e representa \textbf{forte} contribuição ao processo educativo. 655 O laboratório de informática pode significar o acesso a \textbf{uma} gama de conhecimentos quase ilimitados, ou \textbf{um} espaço de distração.
Cabe ao professor ter clareza dos objetivos pedagógicos da \textbf{disciplina} para conduzir, de forma satisfatória, \textbf{uma} aula com o uso de computador.
Desenvolver no aluno a capacidade de pesquisar na internet, ensiná -lo a separar fontes de informações confiáveis de outras de origem duvidosa ou desprovidas de fundamentação, apresenta -se como \textbf{um} desafio permanente. 656 A utilização de filmes e de documentários tem se \textbf{revelado} num \textbf{forte} aliado didático dos professores. \textbf{Contudo}, este recurso deve ser empregado de forma \textbf{bastante} criteriosa, portanto faz -se necessário \textbf{que} o filme apresentado, além de estar relacionado com o conteúdo da \textbf{disciplina} a ser trabalhado naquele período, faça parte do planejamento do docente e possibilite, através do apelo sensível, reflexões, questionamentos e debates.
Pode se \textbf{revelar} como \textbf{uma} oportunidade \textbf{interessante} para a produção de breves dissertações, como propostas para reflexão e análises sobre os conteúdos. 657 Enfim, é importante explorar a incontestável força do audiovisual na nossa cultura. \textbf{Uma} boa biblioteca constitui -se num importantíssimo recurso \textbf{didático-pedagógico} para \textbf{uma} escola \textbf{comprometida} com \textbf{uma} educação de qualidade, porém a biblioteca deve ser visitada, explorada e devidamente aproveitada como espaço fundamental.
A biblioteca oferece \textbf{uma} vasta gama de conhecimentos e de recursos pedagógicos com possibilidades de leituras ilimitadas.
O primeiro desafio \textbf{que} se coloca aos professores é o incentivo à leitura e à escrita, associados ao estudo e à exploração dos recursos pedagógicos disponíveis.
É também importante incutir nas alunas/os o zelo e o cuidado no manuseio dos livros, revistas e periódicos disponibilizados pela biblioteca, bem como a atitude de respeito e silêncio \textbf{que} deve prevalecer num espaço destinado ao estudo.
Tudo isto \textbf{exige} do professor a preparação dos alunas/os para as atividades desenvolvidas na biblioteca. 658 A \textbf{disciplina} de Sociologia reveste -se de \textbf{uma} condição privilegiada, no sentido de poder contribuir para \textbf{que} espaços institucionais escolares sejam fortalecidos, dinamizados ou mesmo criados.
Vale destacar o papel do Grêmio Escolar.
Os professores da \textbf{disciplina} podem e devem contribuir para \textbf{que} os alunas/os compreendam \textbf{que} democracia se constrói com participação efetiva, \textbf{que} cidadania se conquista com esforço e \textbf{que} \textbf{exige} \textbf{disciplina} e responsabilidade. \textbf{Urge} recuperar, nas novas gerações, o interesse pela vida política, não apenas restrita a \textbf{uma} concepção político-partidária, mas como exercício de \textbf{uma} cidadania ativa \textbf{que} considere as aspirações e necessidades coletivas. 659 A instituição ou fortalecimento do Grêmio Escolar emerge como \textbf{uma} possibilidade de \textbf{despertar} o corpo \textbf{discente} para participar e contribuir, de forma efetiva, no enfrentamento dos problemas da escola.
Além disso, pode contribuir na articulação com a comunidade em prol do fortalecimento do Conselho Escolar.
As políticas públicas no campo da educação têm apontado na \textbf{direção} de \textbf{uma} escola \textbf{comprometida} com os valores democráticos, especialmente, a participação e o exercício da cidadania.
Neste sentido, o envolvimento das lideranças estudantis no fortalecimento do Conselho Escolar, além de assegurar -lhe a legitimidade, reveste -se numa oportunidade de aprendizado político. 660 A \textbf{disciplina} de Sociologia tem grandes possibilidades de ampliar o seu espaço de atuação, inclusive extrapolando os muros escolares.
O professor pode propor o desenvolvimento de algum projeto social na escola, engajando os alunas/os, ou articular com outras instituições, iniciativas a serem desenvolvidas na comunidade do entorno da escola. 661 Percebe -se assim \textbf{que} há múltiplas formas e espaços para se ensinar Sociologia, entretanto, a orquestração afinada desta diversidade constitui \textbf{um} desafio aos professores. 662 A \textbf{disciplina} de Sociologia, ao tratar do \textbf{homem} numa perspectiva social, assume como objeto, o estudo das instituições, dos problemas sociais e do comportamento humano.
Esta condição facilita \textbf{um} ensino contextualizado e interdisciplinar.
Trabalhar temas pertinentes à realidade social pode \textbf{despertar} grande interesse no \textbf{educando}, criando oportunidades para \textbf{que} ele faça \textbf{uma} leitura fundamentada em conceitos e \textbf{teorias} próprias do conhecimento sociológico. 663 Para Ileizi Silva ( 2004 ), a problemática sobre o perfil dos professores de Sociologia \textbf{que} atuam nas salas de aula na \textbf{disciplina} Sociologia, mas \textbf{que} não possuem formação específica na Licenciatura em Ciências Sociais pode comprometer o ensino dessa \textbf{disciplina} no ensino médio, \textbf{visto} a ausência de conhecimentos \textbf{metodológicos} e de \textbf{abordagem} sobre o tratamento didático do saber sociológico em saber escolar por parte do professor \textbf{que} ministra a aula.
Isso pode dificultar o entendimento por parte dos alunas/os sobre as grandes questões sociais \textbf{que} estão presentes no cotidiano da vida em sociedade, e \textbf{que} fazem parte do debate das Ciências Sociais e, em especial, da Sociologia.
São questões e problemas sociais \textbf{que} precisam ser analisados a partir dos domínios das Ciências Sociais/Sociologia, pois quando se analisam quaisquer temáticas desse tipo por profissionais de outras ciências \textbf{que} desconhecem o campo \textbf{metodológico} do ensino de Sociologia para a educação básica ( Estranhamento e Desnaturalização ) - mesmo sendo profissionais da área de ciências humanas – pode -se fugir do foco de análise específico, distante das especificidades dos conhecimentos das Ciências Sociais.
A tendência \textbf{aqui} poderá ser \textbf{uma} tentativa de análise das questões sociais, mas sem o aprofundamento \textbf{teórico} necessário.
A análise passa a ter fragilidades de explicação e compreensão sociológicas. 664 O ensino de Sociologia pressupõe, por parte do professor, o \textbf{compromisso} de desenvolver a mediação pedagógica, desenvolvendo estratégias \textbf{que} possibilitem adaptar os conteúdos ao nível intelectual e cognitivo do corpo \textbf{discente}, constituído, em sua maioria, de jovens. \textbf{Um} pressuposto fundamental reside em reconhecer \textbf{que} o objetivo do ensino de Sociologia não é o de formar sociólogos, mas permitir \textbf{que} o aluno tenha contato com conceitos e \textbf{teorias} \textbf{que} permitam \textbf{uma} compreensão mais racional da realidade, possibilitando \textbf{que} ele se situe \textbf{enquanto} indivíduo.
Elaborar esta mediação pedagógica parece \textbf{uma} \textbf{tarefa} \textbf{simples}, entretanto, \textbf{exige} sensibilidade, perspicácia e \textbf{um} bom domínio das \textbf{teorias} sociológicas, bem como alguma experiência de sala de aula.
O exercício da leitura, da interpretação, da reflexão e do debate deve se constituir no processo inicial do aprendizado de Sociologia, além do \textbf{estímulo} à produção de textos dissertativos sobre os temas e problemas depurados através do estudo e da discussão. \textbf{Um} desafio importante para o professor reside em propor \textbf{tarefas} \textbf{que} não \textbf{exijam} apenas o recurso intelectual da memória, sob pena de condenar a \textbf{disciplina} ao estigma de “ matéria decorativa ”, usualmente atribuído às \textbf{disciplinas} da área de ciências humanas. 665 A Sociologia apresenta complexidade e abrangência \textbf{que} a torna \textbf{um} conhecimento vocacionado para o exercício da compreensão, da discussão, da análise e da crítica.
Algumas práticas pedagógicas podem significar exercícios \textbf{que} possam ser aproveitados como processo avaliativo.
A título de ilustração, sugere -se a leitura e a interpretação sociológica de textos ( inclusive de revistas semanais, jornais e outros gêneros ), estudos dirigidos, seminários sobre temas sociais relevantes ( violência, desemprego, pobreza, prostituição, \textbf{exclusão} social e outros ), visitas a museus e instituições públicas ( Câmaras Municipais, Assembleias Legislativas, Centros Administrativos, dentre outros ), bem como a Organizações Não-Governamentais ( ONG ) existentes na comunidade, estimulando a articulação entre o aprendizado dos conceitos e das \textbf{teorias} e a compreensão da realidade social. 666 \textbf{Certamente}, o maior desafio do professor de Sociologia relaciona -se à possibilidade de \textbf{despertar} no \textbf{educando} o interesse pela pesquisa.
Cabe ao professor a \textbf{tarefa} de orientar o \textbf{que} significa \textbf{uma} pesquisa, quais procedimentos \textbf{metodológicos} devem ser adotados, estimulando a ideia de \textbf{que} produzir conhecimento científico pressupõe a articulação entre conceitos, \textbf{teoria} e realidade e \textbf{que} o produto final da pesquisa é, de certo modo, \textbf{uma} reconstrução elaborada a partir das contribuições de outros \textbf{pesquisadores}.
Nunca será pouco advertir quanto aos aspectos ético-legais das pesquisas, levando o aluno a reconhecer \textbf{que}, se apropriar indevidamente do conhecimento alheio, copiar sem mencionar a fonte, além de antiético é ilegal, pois fere o princípio da propriedade intelectual. 667 A diversidade e multiplicidade sobre o ensino da sociologia tem se potencializado e merecido \textbf{um} debate crescente, através dos cursos de licenciatura, congressos, encontros científicos, seminários, programas universitários, sindicatos e associações da área, embora com tantos velhos desafios postos no cenário nacional, a ciência se \textbf{renova} e se fortalece cada vez mais, assim como o tempo, implacável, segue \textbf{adiante}.
A Sociologia segue.

% matched lemmas: abordagem, adiante, aproximar, aqui, basear, bastante, categoria, central, certamente, comprometido, compromisso, conseguir, contemporaneidade, contemporâneo, contudo, desempenhar, desencadear, despertar, dialético, didático-pedagógico, direção, discente, disciplina, dizer, docência, educar, enquanto, estímulo, exclusão, exigir, expressivo, forte, fragmentação, historicamente, homem, humanidade, imbricar, imperativo, implicar, influência, interessante, marca, metodológico, moderno, método, necessitar, ora, perder, pesquisador, que, raciocínio, relegar, remeter, renovar, revelar, simples, sujeito, século, talvez, tarefa, teoria, teórico, um, urgir, ver, viver
\end{textsample}
